\chapter{عدالت انتقالی در دوران گذار}
\label{app:justice}

در این بخش، مدل‌های تفصیلی عدالت انتقالی، مقایسه‌ی رویکردهای جهانی و نقشه راه اجرایی برای ایران ارائه می‌شود.

\section{مدل‌های عدالت انتقالی: ۴ ستون و ۵ ستون}

رویکرد استاندارد بین‌المللی بر اساس ۴ ستون اصلی است که در طرح مهستان برای جامعیت بیشتر به ۵ ستون ارتقا یافته است.

\begin{modernbox}{ستون‌های عدالت انتقالی}
\begin{enumerate}
    \item \textbf{حقیقت‌یابی:} مستندسازی سیستماتیک نقض حقوق بشر.
    \item \textbf{عدالت کیفری:} محاکمه‌ی عاملان اصلی جنایات.
    \item \textbf{جبران خسارت:} ترمیم آسیب‌های مادی و معنوی قربانیان.
    \item \textbf{اصلاحات نهادی:} تغییر ساختارهایی که سرکوب را ممکن کردند.
    \item \textbf{آشتی و حافظه:} (ستون پنجم) تلاش برای التیام زخم‌های اجتماعی و جلوگیری از فراموشی.
\end{enumerate}
\end{modernbox}

\section{مقایسه‌ی رویکردها به عاملان جنایت}

یکی از حساس‌ترین تصمیمات دوران گذار، نحوه‌ی برخورد با عاملان در سطوح مختلف است.

\begin{table}[htbp]
\centering
\caption{طبقه‌بندی عاملان و رویکرد پیشنهادی}
\begin{tabular}{R{4cm} L{6cm} R{4cm}}
\hline
\textbf{سطح مسئولیت} & \textbf{شرح و مصادیق} & \textbf{رویکرد پیشنهادی} \\
\hline
رهبران و فرماندهان عالی & اعضای شورای عالی، فرماندهان ارشد سپاه و اطلاعات، قضات صادرکننده احکام اعدام سیاسی & محاکمه کامل و علنی، بدون امکان عفو \\
فرماندهان میانی و مجریان ارشد & فرماندهان استانی، بازجویان ارشد، دادستان‌های فعال در سرکوب & محاکمه با امکان تخفیف در صورت همکاری کامل \\
مجریان رده‌پایین & سربازان و مأموران عادی، افراد تحت فشار یا تهدید & عفو مشروط در ازای اعتراف کامل و عذرخواهی \\
همکاران غیرمستقیم & مقامات دولتی غیرامنیتی، رسانه‌های تبلیغاتی، روحانیون حامی سرکوب & حقیقت‌یابی، شفافیت مالی و محدودیت‌های شغلی موقت \\
\hline
\end{tabular}
\end{table}

\newpage
\begin{landscape}
\section{خط زمانی اجرایی عدالت انتقالی}

\begin{pgfgantt}[
    vgrid, hgrid,
    title/.style={fill=MainBlue, text=white},
    canvas/.style={fill=white},
    group/.style={fill=SlateGray},
    bar/.style={fill=LightBlue},
    milestone/.style={fill=AccentGold},
    y unit chart=0.7cm,
    x unit=0.5cm,
    time slot format=isodate
]{2025-10-01}{2027-10-01}
    \gantttitle{تقویم اجرایی عدالت انتقالی (۲۰۲۵-۲۰۲۷)}{25} \\
    \gantttitlecalendar{year, month=shortname} \\
    
    \ganttgroup{حقیقت‌یابی}{2025-10-01}{2027-06-01} \\
    \ganttbar{تشکیل کمیسیون}{2025-10-01}{2025-11-30} \\
    \ganttbar{ثبت شهادت قربانیان}{2025-12-01}{2026-11-30} \\
    \ganttbar{گزارش نهایی جامع}{2027-03-01}{2027-06-01} \\

    \ganttgroup{عدالت کیفری}{2025-10-01}{2027-09-01} \\
    \ganttbar{بازداشت عاملان اصلی}{2025-10-01}{2025-12-31} \\
    \ganttbar{محاکمه فرماندهان عالی}{2026-06-01}{2027-06-01} \\
    
    \ganttgroup{جبران خسارت}{2025-10-01}{2027-10-01} \\
    \ganttmilestone{آزادی زندانیان سیاسی}{2025-10-01} \\
    \ganttbar{پرداخت‌های فوری}{2025-11-01}{2026-01-31} \\
    \ganttbar{برنامه ملی غرامت}{2026-05-01}{2027-10-01} \\

    \ganttgroup{آشتی و حافظه}{2026-01-01}{2027-10-01} \\
    \ganttbar{طراحی موزه ملی حافظه}{2026-06-01}{2027-06-01} \\
    \ganttmilestone{اعلام روز ملی یادبود}{2026-10-01} \\
    \ganttmilestone{افتتاح موزه}{2027-10-01} 
\end{pgfgantt}
\end{landscape}

\section{عدالت ترمیمی در مقابل عدالت تنبیهی}

{\centering
\begin{tikzpicture}[node distance=2.5cm, auto, scale=0.8, every node/.style={scale=0.8}]
    \tikzstyle{block} = [rectangle, draw, fill=blue!10, text width=6cm, text centered, rounded corners, minimum height=3.5em, font=\small]
    \tikzstyle{line} = [draw, -Latex, line width=1pt]

    \node [block, fill=SoftRed] (ret) {\textbf{عدالت تنبیهی} \\ تمرکز بر مجازات مجرم \\ سؤال: چه قانونی نقض شد؟};
    \node [block, fill=SoftGreen, below=1.5cm of ret] (res) {\textbf{عدالت ترمیمی} \\ تمرکز بر ترمیم آسیب \\ سؤال: چه کسی آسیب دید؟};
    \node [block, fill=LightBlue, below=1.5cm of res] (trans) {\textbf{عدالت تحول‌گرا} \\ تمرکز بر تغییر ساختاری \\ سؤال: چه ساختاری عامل جنایت بود؟};

    \path [line] (ret) -- (res);
    \path [line] (res) -- (trans);

    \node [right=1cm of res, text width=6cm, fill=SoftGold, draw, rounded corners, align=center] (hybrid) {
        \textbf{رویکرد ترکیبی مهستان:} \\
        ۱. تنبیهی برای سران سرکوب \\
        ۲. ترمیمی برای بدنه و جامعه \\
        ۳. تحول‌گرا برای اصلاح نهادها
    };
\end{tikzpicture}
\par}
